\chapter{Descrizione delle funzioni}

In questa sezione vengono descritte le funzioni principali di ciascun file del backend, con spiegazioni dettagliate e riferimenti al codice.

\section{main.rs}
Le funzioni principali sono:

\begin{lstlisting}[language=Rust, caption={Funzione main}]
#[tokio::main]
async fn main() {
    // Avvio del server e inizializzazione
}
\end{lstlisting}

\section{router.rs}
Definisce la funzione che costruisce il router dell'applicazione:

\begin{lstlisting}[language=Rust, caption={Costruzione del router}]
pub fn create_router(state: AppState) -> Router {
    // Definizione delle rotte e collegamento ai controller
}
\end{lstlisting}

\section{controllers/}
Ogni controller espone funzioni per gestire le richieste HTTP/WS. Ad esempio:

\begin{lstlisting}[language=Rust, caption={Esempio handler in authentication\_controller.rs}]
pub async fn login_handler(Json(payload): Json<LoginDTO>, State(state): State<AppState>) -> Result<Json<UserDTO>, AppError> {
    // Logica di autenticazione
}
\end{lstlisting}

\section{services/}
Le funzioni dei service implementano la logica di business, ad esempio:

\begin{lstlisting}[language=Rust, caption={Esempio funzione in auth\_service.rs}]
pub async fn authenticate_user(db: &PgPool, username: &str, password: &str) -> Result<User, AuthError> {
    // Verifica credenziali
}
\end{lstlisting}

\section{repositories/}
Le funzioni dei repository si occupano di query e accesso ai dati:

\begin{lstlisting}[language=Rust, caption={Esempio funzione in user\_repository.rs}]
pub async fn find_user_by_username(db: &PgPool, username: &str) -> Result<Option<User>, sqlx::Error> {
    // Query al database
}
\end{lstlisting}

\section{DTOs/}
Le funzioni sono principalmente costruttori e conversioni tra tipi.